\section{Supporting Context Variables}\labeledit{sec:cv}

\citet{hu2024layered} present a layered modal type theory which supports pattern
matching on code and establish a normalization proof via a presheaf model. %
However, this form of intentional analysis is not in the most desired form: we cannot
perform recursion on the structure of code. %
This limitations comes from two aspects:
\begin{itemize}
\item STLC lacks a generic notion of types. %
  When we do recursion on, for example, a $\lambda$ expression, the type of the
  $\lambda$ and the type of the body necessarily differ. %
  Therefore, we are not able to formulate a recursion principle without type variables
  as in System F, or a variable of type \texttt{Set} in dependent type theory. %
  This problem naturally goes away if we employ a stronger core type system, so it is
  not our primary concern.
  
\item However, System F or dependent types does not give us a way to capture local
  contexts using a variable. %
  Consider again the $\lambda$ case, even if we have type variables, the recursion on
  the body is in an extended local context, so for the recursion to work, we must be
  able to capture \emph{context variables}, which capture local contexts as global
  variables. 
\end{itemize}

In this section, we focus on context variables. %
We develop the syntactic theory for our $2$-layered contextual model type theory with
context variables, and show its consistency via a reducibility predicate argument. %

\subsection{Well-formedness of Contexts and Types}\labeledit{sec:cv:ctx-types}

With context variables, the type theory becomes ``slightly'' dependently typed, in
that both global and local contexts, and types can depend on context variables, so
their well-formedness requires dedicated judgments. %
The syntax of contexts and types is:
\begin{alignat*}{2}
  i & && \tag{Layer, $i \in[0, 1]$} \\
  x, y & && \tag{Local variables} \\
  u & && \tag{Global variables} \\
  g & && \tag{Context variables} \\
  S, T &:=&&\ \Nat \sep \cont T \sep S \func T \sep (g : \Ctx) \STo T
             \tag{Types, \Typ} \\
  B &:=&&\ u : (\judge T) \sep g : \Ctx
          \tag{Global bindings} \\
  \Phi, \Psi &:= &&\ \cdot \sep \Phi, B
  \tag{Global contexts} \\
  \Gamma, \Delta &:= &&\ \cdot \sep g \sep \Gamma, x : T
                        \tag{Local contexts}
\end{alignat*}
Their well-formedness judgments are:
\begin{mathpar}
  \inferrule
  { }
  {\vdash \cdot}

  \inferrule
  {\vdash \Psi}
  {\vdash \Psi, g : \Ctx}
  
  \inferrule
  {\ljudge 0 \Gamma \\ \ljudge 0 T}
  {\vdash \Psi, u : (\judge T)}

  \inferrule
  {\vdash \Psi}
  {\ljudge i \cdot}

  \inferrule
  {\vdash \Psi \\ g : \Ctx \in \Psi}
  {\ljudge i{g}}

  \inferrule
  {\ljudge i \Gamma \\ \ljudge i T}
  {\ljudge[\Psi] i{\Gamma, x : T}}

  \inferrule
  {\vdash \Psi}
  {\ljudge i \Nat}

  \inferrule
  {\ljudge i S \\ \ljudge i T}
  {\ljudge i {S \func T}}

  \inferrule
  {\ljudge 0 \Delta \\ \ljudge 0 T}
  {\ljudge 1 {\cont[\Delta] T}}

  \inferrule
  {\ljudge[\Psi, g : \Ctx] 1 T}
  {\ljudge 1 {(g : \Ctx) \STo T}}
\end{mathpar}
$\vdash \Psi$ states the well-formedness of a global context $\Psi$. %
There are two kinds of bindings in a global context: either $x : (\judge T)$ as in 
previous modal type theory, or $g : \Ctx$ which is a context variable representing a local
context. %
Eventually, $g$ will be substituted by a concrete and well-formed local context. %
Due to the introduction of context variables, the system clearly becomes
dependently typed. %
$\ljudge i \Gamma$ states the well-formedness of a local context $\Gamma$ at
layer $i$. %
Since we are dealing with a $2$-layered system now, we know $i \in [0, 1]$. %
The base case of a local context can either be an empty local context, or be
a well-scoped context variable. %
The layer $i$ is propagated to the well-formedness judgment of types $\ljudge i T$,
which states that $T$ is well-formed in $\Psi$ at layer $i$. %
Notice that the well-formedness of $T$ does not depend on any local context. %
This judgment essentially combines $\iscore T$ and $\istype T$
predicates in \citet{hu2024layered}, with context variables taken into consideration. %
Context variables is introduced by a special meta-function type $(g : \Ctx)
\STo T$, which pushes a context variable $g$ to the global context. %
This type can be seen as a ``meta-function space'' in which we define macros or
meta-programs, so intuitively, this type and its terms can only live at layer $1$. %
We prove the presupposition lemma of these three judgments:
\begin{lemma}[Presupposition]$ $
  \begin{itemize}
  \item If $\ljudge i \Gamma$,  then $\vdash \Psi$.
  \item If $\ljudge i T$, then $\vdash \Psi$.
  \end{itemize}
\end{lemma}
\begin{lemma}[Lifting] $ $
  \begin{itemize}
  \item If $\ljudge 0 \Gamma$,  then $\ljudge 1 \Gamma$.
  \item If $\ljudge 0 T$,  then $\ljudge 1 T$.
  \end{itemize}
\end{lemma}

\subsection{Weakenings}

Similar to previous layered modal type theories, we need two notions of weakenings: a global one and
a local one. %
In this section, due to context variables, we change the definition of weakenings
based on their counterparts in \citet{hu2024layered}:
\begin{align*}
  \gamma &:= \id \sep q(\gamma) \sep p(\gamma)
  \tag{Global weakenings}
  \\
  \tau &:= \id \sep q(\tau) \sep p(\tau)
  \tag{Local weakenings}
\end{align*}
We will use a global weakening $p(\id)$ in the typing rule of meta-functions in
$\Gamma[p(\id)]$ to account for the insertion of $g: \Ctx$ to the global context
$\Psi$. %
In the following, we examine the properties of weakenings. %
First, we define the composition of weakenings. %
We only define the one for global weakenings and the one for local weakenings is
completely identical:
\begin{align*}
  \id \circ \gamma' &:= \gamma' \\
  \gamma \circ \id &:= \gamma \\
  p(\gamma) \circ q(\gamma') &:= p(\gamma \circ \gamma') \\
  q(\gamma) \circ q(\gamma') &:= q(\gamma \circ \gamma') \\
  \gamma \circ p(\gamma') &:= p(\gamma \circ \gamma')
\end{align*}
\begin{lemma}[Associativity] $ $
  \begin{itemize}
  \item $(\gamma \circ \gamma') \circ \gamma'' = \gamma \circ (\gamma' \circ \gamma'')$
  \item $(\tau \circ \tau') \circ \tau'' = \tau \circ (\tau' \circ \tau'')$
  \end{itemize}
\end{lemma}
Then we apply global weakenings to types and contexts:
\begin{align*}
  \Nat[\gamma] &:= \Nat \\
  S \func T[\gamma] &:= (S[\gamma]) \func (T[\gamma]) \\
  \cont T [\gamma] &:= \cont[\Gamma[\gamma]]{T[\gamma]} \\
  (g: \Ctx) \STo T[\gamma] &:= (g: \Ctx) \STo (T[q(\gamma)]) \\[0.25em]
  \cdot[\gamma] &:= \cdot \\
  g[\gamma] &:= g
                 \tag{properly weakened depending on the name representation} \\
  \Gamma, x : T[\gamma] &:= (\Gamma[\gamma]), x : (T[\gamma])
\end{align*}
This definition of global weakenings admits the following lemma:
\begin{lemma}[Algebra of global weakenings] $ $
  \begin{itemize}
  \item $T[\id] = T$
  \item $\Gamma[\id] = \Gamma$
  \item $T[\gamma][\gamma'] = T[\gamma \circ \gamma']$
  \item $\Gamma[\gamma][\gamma'] = \Gamma[\gamma \circ \gamma']$
  \end{itemize}
\end{lemma}
The well-formedness of weakenings is given by the following rules:
\begin{mathpar}
  \inferrule
  {\vdash \Psi}
  {\id: \Psi \To_g \Psi}

  \inferrule
  {\gamma : \Psi \To_g \Phi \\ \judge[\Psi] B}
  {p(\gamma) : \Psi, B \To_g \Phi}

  \inferrule
  {\gamma : \Psi \To_g \Phi \\ \judge[\Phi] B \\ \judge[\Psi]{B[\gamma]}}
  {q(\gamma) : \Psi, B[\gamma] \To_g \Phi, B}
  \\
  
  \inferrule
  {\ljudge[\Psi] i \Gamma}
  {\id: \Psi;\Gamma \To_i \Gamma}

  \inferrule
  {\tau : \Psi; \Gamma \To_i \Delta \\ \ljudge i T}
  {p(\tau) : \Psi; \Gamma, x : T \To_i \Delta}

  \inferrule
  {\tau : \Psi; \Gamma \To_i \Delta \\ \ljudge i T}
  {q(\tau) : \Psi; \Gamma, x : T \To_i \Delta, x : T}
\end{mathpar}
where the well-formedness of global bindings $\judge[\Psi] B$ is given as follows:
\begin{mathpar}
  \inferrule
  {\vdash \Psi}
  {\judge[\Psi]{g: \Ctx}}

  \inferrule
  {\ljudge 0 \Gamma \\ \ljudge 0 T}
  {\judge[\Psi]{u: (\judge T)}}
\end{mathpar}
The identity case for local weakenings is slightly more complex because we must take
the context variables into consideration. %

Then we can prove the following global weakening lemma:
\begin{lemma}[Global weakenings]\labeledit{lem:cvar:glob-wk-ty} $ $
  \begin{itemize}
  \item If $\ljudge[\Phi] i \Gamma$ and $\gamma : \Psi \To_g \Phi$, then $\ljudge i{ \Gamma[\gamma]}$.
  \item If $\ljudge[\Phi] i T$ and $\gamma : \Psi \To_g \Phi$, then $\ljudge i{ T[\gamma]}$.
  \end{itemize}
\end{lemma}
\begin{proof}
  Mutual induction on $\ljudge[\Phi] i \Gamma$ and $\ljudge[\Phi] i T$.
\end{proof}

The action of global weakenings on a type at layer $0$ is no-op:
\begin{lemma}
  If $\ljudge 0 T$, then $T[\gamma] = T$.
\end{lemma}
Global weakenings do not really affect local weakenings:
\begin{lemma}
  If $\gamma : \Psi \To_g \Phi$ and $\tau : \Phi; \Gamma \To_i \Delta$, then $\tau :
  \Psi; \Gamma[\gamma] \To_i \Delta[\gamma]$. 
\end{lemma}
\begin{proof}
  Induction on $\tau : \Phi; \Gamma \To_i \Delta$.
\end{proof}

The actions of local weakenings only affect terms so they will be looked into in the
next section, after we consider the syntax of the type theory.

\subsection{Syntax and Typing}

In this section, we define the syntax of the type theory with context variables. %
To isolate concerns, we use $\tletbox$ for elimination, instead of pattern matching. %
Nevertheless, pattern matching on code should work with proper adjustments to our
development:
\begin{alignat*}{2}
  m & && \tag{Natural numbers, $\N$} \\
  \delta &:=&& \cdot_{g?}^{m} \sep \wk_g^m \sep \delta, t/x \tag{Local substitutions} \\
  s, t &:=&&\ x \sep u^\delta \tag{Terms, $\Exp$} \\
  & && \sep \ze \sep \su t
  \tag{natural numbers}\\
  & && \sep \boxit t \sep \letbox u s t
  \tag{box}\\
  & &&\sep \lambda x. t \sep s\ t \tag{functions}  \\
  & &&\sep \Lambda g. t \sep t ~\$~ \Gamma \tag{meta-functions} 
\end{alignat*}
Similar to before, we have natural numbers as our base type. %
To construct and eliminate meta-functions, we have $\Lambda g. t$ and $t ~\$~ \Gamma$
respectively. %
Since we have contextual types, each global variable must be associated with a local
substitution. %
In a local substitution, due to how a local context is structured, there are also two
base cases: it can either be an empty local substitution $\cdot_{g?}^{m}$, or a weakening $\wk^m_g$ of a
context variable $g$. %
The number $m$ associated with both cases are the number of local weakening $p$'s. %
Effectively, $m$ equals to the length of the codomain local context in the typing
judgment, as we will specify later in this section. %
This number can be fetched from a local substitution by the following function:
\begin{align*}
  \widehat{\cdot_{g?}^m} &:= m \\
  \widehat{\wk_{g}^m} &:= m \\
  \widehat{\delta, t/x} &:= \widehat{\delta}
\end{align*}
In addition, $\cdot$ is optionally associated with a context variable $g$. %
If there is such a $g$, it is the base case of the codomain local context. %
These information are necessary in order to define the local and global substitution
operations. %
We rely on the following function to return the context variable inside of a local
substitution, if it exists:
\begin{align*}
  \widecheck{\cdot_{g?}^m} &:= g? \\
  \widecheck{\wk_g^m} &:= g \\
  \widecheck{\delta, t/x} &:= \widecheck{\delta}
\end{align*}
Next, we define the action of global weakenings on terms and local substitutions:
\begin{align*}
  x[\gamma] &:= x \\
  u^\delta[\gamma] &:= u^{\delta[\gamma]}
                     \tag{with $u$ properly weakened depending on name representation} \\
  \ze[\gamma] &:= \ze \\
  \su t[\gamma] &:= \su {(t[\gamma])} \\
  \boxit t[\gamma] &:= \boxit {(t[\gamma])} \\
  \letbox u s t[\gamma] &:= \letbox u {s[\gamma]} {(t[q(\gamma)])} \\
  \lambda x. t[\gamma] &:= \lambda x. {(t[\gamma])} \\
  t~s[\gamma] &:= (t[\gamma])~(s[\gamma]) \\
  \Lambda g. t[\gamma] &:= \Lambda g. {(t[q(\gamma)])} \\
  t~\$~\Gamma[\gamma] &:= (t[\gamma])~\$~(\Gamma[\gamma]) \\[0.25em]
  \wk_g^m[\gamma] &:= \wk_g^m
                     \tag{with $g$ properly weakened} \\
  \cdot_{g?}^m[\gamma] &:= \cdot_{g?}^m
                         \tag{with $g$ properly weakened if exists} \\
  \delta, t/x[\gamma] &:= (\delta[\gamma]), (t[\gamma])/x
\end{align*}
The application of global weakenings satisfy the following lemma:
\begin{lemma}[Algebra of global weakenings]$ $
  \begin{itemize}
  \item $t[\id] = t$
  \item $\delta[\id] = \delta$
  \item $t[\gamma][\gamma'] = t[\gamma \circ \gamma']$
  \item $\delta[\gamma][\gamma'] = \delta[\gamma \circ \gamma']$
  \end{itemize}
\end{lemma}

Applying local weakenings on terms and local substitutions is defined as follows:
\begin{align*}
  x[\tau] &:= x
  \tag{properly weakened} \\
  u^\delta[\tau] &:= u^{\delta[\tau]} \\
  \ze[\tau] &:= \ze \\
  \su t[\tau] &:= \su {(t[\tau])} \\
  \boxit t[\tau] &:= \boxit t \\
  \letbox u s t[\tau] &:= \letbox u {s[\tau]} {(t[\tau])} \\
  \lambda x. t[\tau] &:= \lambda x. {(t[q(\tau)])} \\
  t~s[\tau] &:= (t[\tau])~(s[\tau]) \\
  \Lambda g. t[\tau] &:= \Lambda g. {(t[\tau])} \\
  t~\$~\Gamma[\tau] &:= (t[\tau])~\$~\Gamma \\[0.25em]
  \wk_g^m[\tau] &:= \wk_g^{m + m'}
  \tag{where $m'$ is the number of $p$ constructor in $\tau$} \\
  \cdot_{g?}^m[\tau] &:= \cdot_{g?}^{m + m'}
                       \tag{where $m'$ is the number of $p$ constructor in $\tau$} \\
  \delta, t/x[\tau] &:= (\delta[\tau]), (t[\tau])/x
\end{align*}
The application of local weakenings satisfy the following lemma:
\begin{lemma}[Algebra of local weakenings]$ $
  \begin{itemize}
  \item $t[\gamma][\tau] = t[\tau][\gamma]$
  \item $\delta[\gamma][\tau] = \delta[\tau][\gamma]$
  \item $t[\tau][\tau'] = t[\tau \circ \tau']$
  \item $\delta[\tau][\tau'] = \delta[\tau \circ \tau']$
  \item $t[\delta][\tau] = t[\delta[\tau]]$
  \item $(\delta \circ \delta')[\tau] = \delta \circ (\delta'[\tau])$
  \end{itemize}
\end{lemma}

Then the weakenings of dual-contexts are defined as just tuples of global and local
weakenings:
\begin{align*}
  &\inferrule
  {\gamma : \Psi \To_g \Phi \\ \tau : \Psi;\Gamma \To_i \Delta[\gamma]}
  {\gamma; \tau : \Psi; \Gamma \To_i \Phi; \Delta} \\
  &t[\gamma; \tau] := t[\gamma][\tau] \\
  &\delta[\gamma; \tau] := \delta[\gamma][\tau]
\end{align*}
To disambiguate, when we apply a local weakening literal, we usually pair it with an
identity global weakening. %
For example, $t [ \id; p(\id)]$ \emph{locally} weakens $t$ with a local weakening $p(\id)$. %
When we write $t[p(\id)]$, we mean that $t$ is \emph{globally} weakened by $p(\id)$. %
The correctness lemma for these operations can only be proved after defining the typing
rules. %
In order to define the typing rules, we must first give the definition of global
substitutions of types and local contexts:
\begin{alignat*}{2}
  \sigma &:=&& \cdot \sep \sigma, t/u \sep \sigma, \Gamma/g
               \tag{Global substitutions}
\end{alignat*}
Global substitutions can obviously be globally weakened by iteratively applying a
global weakening to the terms and local contexts within. %
The global substitution operation of types and local contexts is defined as follows:
\begin{align*}
  \Nat [\sigma] &:= \Nat \\
  S \func T [\sigma] &:= (S[\sigma]) \func (T[\sigma]) \\
  \cont T [\sigma] &:= \cont[\Gamma[\sigma]]{T[\sigma]} \\
  (g: \Ctx) \STo T [\sigma] &:= (g: \Ctx) \STo (T [\sigma[p(\id)], g/g])
  \\[0.25em]
  \cdot [\sigma] &:= \cdot \\
  g [\sigma] &:= \sigma(g)
                  \tag{lookup $g$ in $\sigma$; undefined if $g$ is not bound or result is not a
                  local context} \\
  \Gamma, x : T[\sigma] &:= (\Gamma[\sigma]) , x : (T[\sigma])
\end{align*}
For consistency of notations, we often write $q(\sigma)$ for $\sigma[p(\id)], g/g$ or
$\sigma[p(\id)], u^\id/u$. %
This notation relates similar operations of weakenings and substitutions. %
The global substitutions on types satisfy the following lemma:
\begin{lemma}[Algebra of Global Substitutions]$ $
  \begin{itemize}
  \item $T[\gamma][\gamma'] = T[\gamma \circ \gamma]$
  \item $\sigma[\gamma][\gamma'] = \sigma[\gamma \circ \gamma]$
  \item $T[\sigma][\gamma] = T[\sigma[\gamma]]$
  \item $\Gamma[\sigma][\gamma] = \Gamma[\sigma[\gamma]]$
  \item $t[\sigma][\gamma] = t[\sigma[\gamma]]$
  \item $\delta[\sigma][\gamma] = \delta[\sigma[\gamma]]$
  \end{itemize}
\end{lemma}

Next, we give the application operation of local substitutions on terms and
composition of local substitutions:
\begin{align*}
  x[\delta] &:= \delta(x)
              \tag{lookup of $x$ in $\delta$} \\
  u^{\delta'}[\delta] &:= u^{\delta' \circ \delta} \\
  \ze[\delta] &:= \ze \\
  \su t[\delta] &:= \su{(t[\delta])} \\
  \lambda x. t[\delta] &:= \lambda x. (t[\delta[\id; p(\id)], x/x]) \\
  t~s[\delta] &:= (t[\delta])~(s[\delta]) \\
  \boxit t [\delta] &:= \boxit t \\
  \letbox u s t[\delta] &:= \letbox u {s[\delta]} {(t[\delta[p(\id)]])} \\
  \Lambda g. t [\delta] &:= \Lambda g. (t[\delta[(p(\id))]]) \\
  t~\$~\Gamma [\delta] &:= (t[\delta]) ~\$~\Gamma \\[.25em]
  \wk_g^m \circ \delta &:= \wk_g^{\widehat\delta} \\
  \cdot^m \circ \delta &:= \cdot_{\widecheck{\delta}}^{\widehat\delta} \\
  \cdot_{g}^m \circ \delta &:= \cdot_{g}^{\widehat\delta} \\
  (\delta', t/x) \circ \delta &:= (\delta' \circ \delta), t[\delta]/x
\end{align*}
Similarly, we might also write $q(\delta)$ for $\delta[\id; p(\id)], x/x$. %
Notice that in the definition of composition, we make use of the $\widehat\delta$
function to fetch the number of weakenings. %
This number is used in the application operation of global substitutions given below
in the application of global substitutions to specify the number of local $p$
weakenings when a context variable is substituted by a concrete context. %
In the composition of $\cdot^m$, we use $\widecheck{\delta}$ to query whether
$\delta$'s codomain context starts from a context variable. %
If it does, then we use that context variable in the result of the composition. %
Local substitutions and global weakenings interact in the following way:
\begin{lemma}[Algebra of Local Substitutions] $ $
  \begin{itemize}
  \item $t[\delta][\gamma] = (t[\gamma])[\delta[\gamma]]$
  \item $(\delta \circ \delta')[\gamma] = (\delta[\gamma]) \circ (\delta'[\gamma])$
  \end{itemize}
\end{lemma}

Now, we define the application of global substitutions to terms and local substitutions:
\begin{align*}
  x[\sigma] &:= x \\
  u^{\delta}[\sigma] &:= \sigma(u)[\delta[\sigma]]
                        \tag{lookup of $u$ in $\sigma$} \\
  \ze[\sigma] &:= \ze \\
  \su t[\sigma] &:= \su{(t[\sigma])} \\
  \lambda x. t[\sigma] &:= \lambda x. (t[\sigma]) \\
  t~s[\sigma] &:= (t[\sigma])~(s[\sigma]) \\
  \boxit t [\sigma] &:= \boxit {(t[\sigma])} \\
  \letbox u s t[\sigma] &:= \letbox u {s[\sigma]} {(t[\sigma[p(\id)], u^\id/u])} \\
  \Lambda g. t [\sigma] &:= \Lambda g. (t[\sigma[(p(\id))], g/g]) \\
  t~\$~\Gamma [\sigma] &:= (t[\sigma]) ~\$~(\Gamma[\sigma]) \\[.25em]
  \wk_g^m [\sigma] &:= \id_{\sigma(g)}[\id; p^m(\id)]
                     \tag{defined only when $\sigma(g)$ is a local context}\\
  \cdot^m [\sigma] &:= \cdot^m \\
  \cdot_{g}^m [\sigma] &:= \cdot^{|\Gamma| + m}
  \tag{if $\sigma(g) = \Gamma$ and $\Gamma$ ends with a $\cdot$} \\
  \cdot_{g}^m [\sigma] &:= \cdot_{g'}^{|\Gamma| + m}
  \tag{if $\sigma(g) = \Gamma$ and $\Gamma$ ends with a $g'$} \\
  (\delta, t/x) [\sigma] &:= (\delta[\sigma]), t[\sigma]/x
\end{align*}
In the definition of global substitution operation, we make use of the local identity
substitution, which is defined through local weakening substitutions as below:
\begin{align*}
  \wk^m_{\cdot} &:=  \cdot^m \\
  \wk^m_{g} &:=  \wk_g^m \\
  \wk^m_{\Gamma, x : T} &:= \wk^{1 + m}_{\Gamma}, x/x
\end{align*}
The local identity substitution is defined as a special case of the local weakening
substitutions by setting $m$ to be $0$:
\begin{align*}
  \id_{\Gamma} := \wk^0_{\Gamma}
\end{align*}
Though we make heavy use of symbol overloading, but hopefully the exact meanings of
symbols should be disambiguated by the surrounding textual contexts. %
Usually, the symbols are designed such that their behaviors remain the same for their
ambiguous readings (e.g. various uses of $\id$ and $\wk$).  %

At last, we need to define the global identity substitution before giving the typing
rules. %
The global identity substitution is defined in the same principle; it is a special
case of the global weakening substitutions:
\begin{align*}
  \wk^m_{\cdot} &:=  \cdot \\
  \wk^m_{\Psi, g : \Ctx} &:=  \wk_{\Psi}^{1+m}, g/g \\
  \wk^m_{\Psi, u : (\judge T)} &:= \wk^{1 + m}_{\Psi}, u^{\id_\Gamma[p^{1+m}(\id)]}/u
\end{align*}
Notice that in the cons case for $u : (\judge T)$, the local identity substitution
$\id_{\Gamma}$ must be weakened by $p^{1+m}(\id)$, because its global typing
environment is weakened by $\wk^{1 + m}_{\Psi}$, which takes the same effect. %
This weakening is necessary to make the typing to go through. %
Then identity is just a special case:
\begin{align*}
  \id_{\Psi} := \wk^0_{\Psi}
\end{align*}
We sometimes omit the subscript for different $\id$, as we know their effects on
types, local contexts, terms and local substitutions are just \emph{identity}. 

The composition of global substitutions is defined intuitively:
\begin{align*}
  \cdot \circ \sigma' &:= \cdot \\
  (\sigma, t/u) \circ \sigma' &:= (\sigma \circ \sigma'), t[\sigma']/u \\
  (\sigma, \Gamma/g) \circ \sigma' &:= (\sigma \circ \sigma'), \Gamma[\sigma']/g
\end{align*}
Essentially, composition just iteratively applies $\sigma'$ to all terms and contexts
in the first global substitution. %
Notice that composition can only be defined, after the applications of global
substitutions to both terms and local contexts are defined. %
This definition does lead to some complication when we tries to prove the global
substitution lemma of terms and local substitutions in the next section. %

At last, the following gives the typing rules of terms and local substitutions.
\begin{mathpar}
  \inferrule
  {\ljudge i \Gamma \\ x : T \in \Gamma}
  {\ltyping{i}{x}{T}}

  \inferrule
  {\ltyping i \delta \Delta \\ u : (\judge[\Delta] T) \in \Psi}
  {\ltyping{i}{u^\delta}{T}}

  \inferrule
  {\ljudge i \Gamma}
  {\ltyping{i}{\ze}{\Nat}}

  \inferrule
  {\ltyping{i}{t}{\Nat}}
  {\ltyping{i}{\su t}{\Nat}}
  
  \inferrule
  {\ltyping[\Psi][\Gamma, x : S]{i}{t}{T}}
  {\ltyping{i}{\lambda x. t}{S \func T}}

  \inferrule
  {\ltyping{i}{t}{S \func T} \\ \ltyping{i}{s}{S}}
  {\ltyping{i}{t\ s}{T}}

  \inferrule
  {\ljudge 1 \Gamma \\ \ltyping[\Psi][\Delta] 0 t T}
  {\ltyping{1}{\boxit t}{\cont[\Delta] T}}
  
  \inferrule
  {\ltyping 1 {s}{\cont[\Delta] T} \\ \mhighlight{\ljudge[\Psi] 0 \Delta} \\
    \mhighlight{\ljudge[\Psi] 0 T} \\
    \ljudge[\Psi] 1{T'} \\
    \ltyping[\Psi, u : (\judge[\Delta] T)][\Gamma[p(\id)]] 1 {t}{T'[p(\id)]}}
  {\ltyping 1 {\letbox u s t} T'}

  \inferrule
  {\ljudge[\Psi] 1 \Gamma \\ \ltyping[\Psi, g: \Ctx][\Gamma[p(\id)]] 1 t T}
  {\ltyping{1}{\Lambda g. t}{(g : \Ctx) \STo T}}

  \inferrule
  {\ltyping{1}{t}{(g : \Ctx) \STo T} \\ \ljudge[\Psi] 0 \Delta}
  {\ltyping{1}{t~\$~\Delta}{T[\id_{\Psi}, \Delta/g]}}
  
  \inferrule
  {\ljudge i \Gamma \\ \text{$\Gamma$ ends with $\cdot$} \\ |\Gamma| = m}
  {\ltyping i {\cdot^m}{\cdot}}

  \inferrule
  {\ljudge i \Gamma \\ g : \Ctx \in \Psi \\ \text{$\Gamma$ ends with $g$} \\ |\Gamma| = m}
  {\ltyping i {\cdot_g^m}{\cdot}}

  \inferrule
  {\ljudge i \Gamma \\ g : \Ctx \in \Psi \\ \text{$\Gamma$ ends with $g$} \\ |\Gamma| = m}
  {\ltyping i {\wk_g^m}{g}}

  \inferrule
  {\ltyping i {\delta}{\Delta} \\ \ltyping i {t}{T}}
  {\ltyping i {\delta, t/x}{\Delta, x : T}}
\end{mathpar}
In the typing rules, there are premises highlighted by \mhighlight{\text{shades}}. %
These shaded premises are necessary to establish the theorem of \emph{presupposition}
or \emph{syntactic validity}. %
After establishing presupposition, these premises can be derived from other premises
and thus technically can be omitted afterwards. %
Moreover, in the rule for global variables, the lookup of a global context $\Psi$ must
consider the effect of global weakenings as follows:
\begin{mathpar}
  \inferrule
  { }
  {u : B[p(\id)] \in \Psi, u : B}

  \inferrule
  {u : B \in \Psi}
  {u : B[p(\id)] \in \Psi , u' : B'}
\end{mathpar}
The typing rules also support lifting:
\begin{lemma}[Lifting] $ $
  \begin{itemize}
  \item If $\ltyping{0}{t}{T}$, then $\ltyping{1}{t}{T}$.
  \item If $\ltyping{0}{\delta}{\Delta}$, then $\ltyping{1}{\delta}{\Delta}$.
  \end{itemize}
\end{lemma}
This lemma ensures that terms at layer $0$ are \emph{included} in layer $1$. 

Typing rules for global substitutions are defined as follows:
\begin{mathpar}
  \inferrule*
  {\vdash \Psi}
  {\typing[\Psi]{\cdot}{\cdot}}

  \inferrule*
  {\typing[\Psi]{\sigma}{\Phi} \\ \ljudge[\Psi] 0 \Gamma \\ \ljudge[\Psi] 0 T \\ \ltyping[\Psi][\Gamma[\sigma]] 0 {t}{T[\sigma]}}
  {\typing[\Psi]{\sigma, t/u}{\Phi, u : (\judge T)}}

  \inferrule*
  {\typing[\Psi]{\sigma}{\Phi} \\ \ljudge 0 \Gamma}
  {\typing[\Psi]{\sigma, \Gamma/g}{\Phi, g: \Ctx}}
\end{mathpar}

In the next section, we establish a set of syntactic properties as a basic sanity
check of the definitions, which are also useful in later section proving normalization. 

\subsection{Syntactic Properties of 2-layered Modal Type Theory with Context Variables}

In this section, we list syntactic properties that eventually leads to the
substitution lemma of terms and local substitutions for global substitutions. %
This lemma is our ``benchmark'' to ensure that the rules for the system make sense. %
During the process, we must establish other necessary syntactic properties. %
We elaborate the proofs an important and selected few. %
Other proofs in this section have been mechanized in Agda.

\begin{lemma}\labeledit{lem:cvar:gsubst-skip} $ $
  \begin{itemize}
  \item If $n$ is the length of $\sigma'$, then $t [q^n(p(id))][\sigma,t/u,\sigma'] =
    t [q^n(p(id))][\sigma,\Gamma/g,\sigma'] = t[\sigma,\sigma']$.
  \item If $n$ is the length of $\sigma'$, then $\delta [q^n(p(id))][\sigma,t/u,\sigma'] =
    \delta [q^n(p(id))][\sigma,\Gamma/g,\sigma'] = \delta[\sigma,\sigma']$.
  \end{itemize}
\end{lemma}
This lemma allows to skip a binding in the middle of a global substitution according
to a global weakening.

A similar lemma holds for local substitutions:
\begin{lemma}\labeledit{lem:cvar:lsubst-skip} $ $
  \begin{itemize}
  \item If $n$ is the length of $\delta'$, then $t [q^n(p(id))][\delta,t/x,\delta'] =
    t[\delta,\delta']$.
  \item If $n$ is the length of $\delta'$, then $\delta [q^n(p(id))] \circ (\delta,t/x,\delta') =
    \delta \circ (\delta,\delta')$.
  \end{itemize}
\end{lemma}


\begin{lemma}[Composition of Global Substitutions]$ $
  \begin{itemize}
  \item $T[\sigma][\sigma'] = T[\sigma \circ \sigma']$
  \item $\Gamma[\sigma][\sigma'] = \Gamma[\sigma \circ \sigma']$
  \end{itemize}
\end{lemma}

\begin{lemma}[Composition and Associativity of Local Substitutions] $ $
  \begin{itemize}
  \item $t[\delta][\delta'] = t[\delta \circ \delta']$
  \item $(\delta \circ \delta') \circ \delta'' = \delta \circ (\delta' \circ \delta'')$
  \end{itemize}
\end{lemma}

\begin{lemma}[Typing of Local Weakening Substitutions]
  If $\ljudge i {\Delta, \Gamma}$, then $\ltyping[\Psi][\Delta, \Gamma] i
  {\wk^{|\Gamma|}_{\Delta}}{\Delta}$. 
\end{lemma}
The corollary is the well-typedness of local identity substitution:
\begin{corollary}
  If $\ljudge i \Gamma$, then $\ltyping i {\id_{\Gamma}}{\Gamma}$. 
\end{corollary}

The next few questions require well-formedness or typing judgments to work:
\begin{lemma}\labeledit{lem:cvar:wk-p-n-ty} $ $
  \begin{itemize}
  \item If $\ljudge i T$, then $T[\wk^n_{\Psi}] = T[p^n(\id)]$. 
  \item If $\ljudge i \Gamma$, then $\Gamma[\wk^n_{\Psi}] = \Gamma[p^n(\id)]$. 
  \end{itemize}
\end{lemma}
This lemma proves that the global weakening substitutions behaves exactly like global
weakenings.

\begin{lemma}[Naturality]$ $
  \begin{itemize}
  \item If $\ljudge[\Psi,g:\Ctx] i T$ and $\gamma : \Phi \To_g \Psi$, then
    $T[\id_{\Psi}, \Gamma/g][\gamma] = T[q(\gamma)][\id_{\Phi},\Gamma[\gamma]/g]$.
  \item If $\ljudge[\Psi,g:\Ctx] i \Delta$ and $\gamma : \Phi \To_g \Psi$, then
    $\Delta[\id_{\Psi}, \Gamma/g][\gamma] = \Delta[q(\gamma)][\id_{\Phi},\Gamma[\gamma]/g]$.
  \end{itemize}
\end{lemma}
The naturality lemma finds correspondence in the characterization of a presheaf
category of a type theory in general, which instructs how $q$ weakenings can be used
to swap a global weakening and a global substitution. %

\begin{lemma}[Local Identity]\labeledit{lem:cvar:loc-id}$ $
  \begin{itemize}
  \item If $\ltyping i t T$, then $t[\id_{\Gamma}] = t$. 
  \item If $\ltyping i \delta \Delta$, then $\delta \circ \id_{\Gamma} = \delta$. 
  \end{itemize}
\end{lemma}
This lemma shows that the local identity substitution has no effect on terms and that
the right identity property of local substitutions. %

Next, we establish the global weakening lemma for typing rules:
\begin{lemma}[Global weakenings]\labeledit{lem:cvar:glob-wk-tm} $ $
  \begin{itemize}
  \item If $\ltyping i t T$ and $\gamma : \Psi' \To_g \Psi$, then
    $\ltyping[\Psi'][\Gamma[\gamma]] i {t[\gamma]}{ T[\gamma]}$.
  \item If $\ltyping i \delta \Delta$ and $\gamma : \Psi' \To_g \Psi$, then
    $\ltyping[\Psi'][\Gamma[\gamma]] i {\delta[\gamma]}{\Delta[\gamma]}$.
  \end{itemize}
\end{lemma}
\begin{proof}
  Mutual induction on $\ltyping[\Phi] i t T$ and $\ltyping[\Phi] i \delta \Delta$. %
  We only consider a few interesting cases:
  \begin{itemize}[label=Case]
  \item
    \begin{mathpar}
      \inferrule
      {\ltyping[\Phi] 1 {s}{\cont[\Delta] T} \\ \ljudge[\Phi] 1 {T'} \\ \ltyping[\Phi, u : (\judge[\Delta] T)][\Gamma[p(\id)]] 1 {t}{T'[p(\id)]}}
      {\ltyping[\Phi] 1 {\letbox u s t} T'}
    \end{mathpar}
    \begin{align*}
      & \ltyping[\Psi][\Gamma[\gamma]] 1 {s[\gamma]}{\cont[\Delta] T[\gamma]}
        \byIH \\
      & \ltyping[\Psi, u : (\judge[\Delta[\gamma]]
        T[\gamma])][\Gamma[p(\id)][q(\gamma)]] 1 {t[q(\gamma)]}{T'[p(\id)][q(\gamma)]}
        \byIH \\
      & \ltyping[\Psi, u : (\judge[\Delta[\gamma]]
        T[\gamma])][\Gamma[p(\gamma)]] 1 {t[q(\gamma)]}{T'[p(\gamma)]}
        \tag{by computation} \\
      & \ltyping[\Psi, u : (\judge[\Delta[\gamma]]
        T[\gamma])][\Gamma[\gamma][p(\id)]] 1 {t[q(\gamma)]}{T'[\gamma][p(\id)]} \\
      & \ltyping[\Psi][\Gamma[\gamma]] 1 {\letbox u s t[\gamma]}{T'[\gamma]}
        \tag{by constructor}
    \end{align*}
    
  \item
    \begin{mathpar}
      \inferrule
      {\ltyping[\Phi, g: \Ctx][\Gamma[p(\id)]] 1 t T}
      {\ltyping[\Phi]{1}{\Lambda g. t}{(g : \Ctx) \STo T}}
    \end{mathpar}
    \begin{align*}
      & q(\gamma) : \Psi', g : \Ctx \To_g \Psi, g : \Ctx
        \tag{by typing rules} \\
      & \ltyping[\Psi', g: \Ctx][\Gamma[p(\id)][q(\gamma)]]
        1{t[q(\gamma)]}{T[q(\gamma)]}
        \byIH \\
      & \ltyping[\Psi', g: \Ctx][\Gamma[p(\id) \circ q(\gamma)]]
        1{t[q(\gamma)]}{T[q(\gamma)]}
        \tag{by algebraic law} \\
      & \ltyping[\Psi', g: \Ctx][\Gamma[p(\gamma)]]
        1{t[q(\gamma)]}{T[q(\gamma)]} \\
      & \ltyping[\Psi', g: \Ctx][\Gamma[\gamma][p(\id)]]
        1{t[q(\gamma)]}{T[q(\gamma)]} \\
      & \ltyping[\Psi'][\Gamma[\gamma]]
        1{\Lambda g. (t[q(\gamma)])}{(g : \Ctx) \STo (T[q(\gamma)])}
        \tag{by typing rule}
    \end{align*}
    
  \item
    \begin{mathpar}
      \inferrule
      {\ltyping[\Phi]{1}{t}{(g : \Ctx) \STo T} \\ \ljudge[\Phi] 0 \Delta}
      {\ltyping[\Phi]{1}{t~\$~\Delta}{T[\id_{\Phi}, \Delta/g]}}
    \end{mathpar}
    \begin{align*}
      & \ltyping[\Psi][\Gamma[\gamma]]{1}{t[\gamma]}{(g : \Ctx) \STo T[\gamma]}
        \byIH \\
      & \ljudge[\Psi] 0 {\Delta[\gamma]}
        \tag{by \Cref{lem:cvar:glob-wk-ty}} \\
      &
        \ltyping[\Psi][\Gamma[\gamma]]{1}{(t[\gamma])~\$~(\Delta[\gamma])}{T[q(\gamma)][\id_{\Psi}, \Delta[\gamma]/g]}
        \tag{by constructor} \\
      & \ltyping[\Psi][\Gamma[\gamma]]{1}{(t[\gamma])~\$~(\Delta[\gamma])}{T[\id_{\Phi},
        \Delta/g][\gamma]}
        \tag{by naturality}
    \end{align*}
    
  \end{itemize}
\end{proof}

\begin{lemma}[Global Weakening]
  If $\typing[\Psi] \sigma \Phi$ and $\gamma : \Psi' \To_g \Psi$, then $\typing[\Psi']{\sigma[\gamma]}\Phi$.
\end{lemma}

\begin{lemma}[Global Weakening Substitutions]
  If $\vdash \Psi,\Phi$, then $\typing[\Psi,\Phi]{\wk^{|\Phi|}_{\Psi}}{\Psi}$. 
\end{lemma}
\begin{proof}
  This lemma is actually requires a bit preliminaries to establish and so this is the
  earliest point where this lemma can be proven. %
  From $\vdash \Psi,\Phi$, we know $\vdash \Psi$, which we do induction on. %
  We consider only one case:
  \begin{mathpar}
    \inferrule
    {\ljudge 0 \Gamma \\ \ljudge 0 T}
    {\vdash \Psi, u : (\judge T)}
  \end{mathpar}
  \begin{align*}
    & \wk^{|\Phi|}_{\Psi, u : (\judge T)} = \wk^{1 + |\Phi|}_{\Psi},
      u^{\id_\Gamma[p^{1+|\Phi|}(\id)]}/u
      \tag{by definition} \\
    & \typing[\Psi, u : (\judge T),\Phi]{\wk^{1 + |\Phi|}_{\Psi}}{\Psi}
      \byIH 
  \end{align*}
  At last, we must prove $\ltyping[\Psi, u : (\judge
  T),\Phi][\Gamma[\wk^{1 + |\Phi|}_{\Psi}]] 0
  {u^{\id_\Gamma[p^{1+|\Phi|}(\id)]}}{T[\wk^{1 + |\Phi|}_{\Psi}]}$. %
  But we know that this goal is the same as the following due to
  \Cref{lem:cvar:wk-p-n-ty}:
  \begin{align*}
    \ltyping[\Psi, u : (\judge
    T),\Phi][\Gamma[p^{1 + |\Phi|}(\id)]] 0
    {u^{\id_\Gamma[p^{1+|\Phi|}(\id)]}}{T[p^{1 + |\Phi|}(\id)]}
  \end{align*}
  It remains to prove that the substitution is well-typed:
  \begin{align*}
    \ltyping[\Psi, u : (\judge
    T),\Phi][\Gamma[p^{1 + |\Phi|}(\id)]] 0
    {\id_\Gamma[p^{1+|\Phi|}(\id)]}{\Gamma[p^{1 + |\Phi|}(\id)]}
  \end{align*}
  This goal is immediate due to \Cref{lem:cvar:loc-id,lem:cvar:glob-wk-tm}.
\end{proof}

\begin{corollary}
  If $\vdash \Psi$, then $\typing[\Psi]{\id_{\Psi}}{\Psi}$. 
\end{corollary}


Finally, we can establish the presupposition of terms and local substitutions:
\begin{lemma}[Presupposition]$ $
  \begin{itemize}
  \item If $\ltyping i t T$, then $\ljudge i \Gamma$ and $\ljudge i T$. 
  \item If $\ltyping i \delta \Delta$, then $\ljudge i \Gamma$ and $\ljudge i \Delta$. 
  \end{itemize}
\end{lemma}
\begin{proof}
  We do a mutual induction.
\end{proof}

Next, we need a similar lemma to \Cref{lem:cvar:wk-p-n-ty} but for terms and local
substitutions:
\begin{lemma}\labeledit{lem:cvar:wk-p-n-tm} $ $
  \begin{itemize}
  \item If $\ltyping i t T$, then $t[\wk^n_{\Psi}] = t[p^n(\id)]$. 
  \item If $\ltyping i \delta \Delta$, then $\delta[\wk^n_{\Psi}] = \delta[p^n(\id)]$. 
  \end{itemize}
\end{lemma}
\begin{proof}
  The proof of this lemma requires an intrigued generalization in order to handle
  extensions of global contexts due to $\tletbox$ and $\Lambda$. %
  Details of the generalization are technical and too elaborate to put in this
  technical report, and thus we choose to leave them in the Agda mechanization for
  readers' reference. 
\end{proof}

Next, we should verify the identity rules of composition of global substitutions. %
Notice that by applying \Cref{lem:cvar:wk-p-n-ty,lem:cvar:wk-p-n-tm}, we obtain the
right identity immediately:
\begin{lemma}
  If $\typing[\Psi] \sigma \Phi$, then $\sigma \circ \id_{\Psi} = \sigma$. 
\end{lemma}

The left identity, on the other hand, requires certain generalization which
must incorporate global weakenings. %
We again leave the details in the Agda mechanization:
\begin{lemma}
  If $\typing[\Psi] \sigma \Phi$, then $\id_{\Phi} \circ \sigma = \sigma$. 
\end{lemma}

Another useful equation is that global substitutions and local weakenings commute:
\begin{lemma}[Commutativity of Global Substitutions and Local Weakenings] $ $
  \begin{itemize}
  \item If $\ltyping i t T$ and $\tau : \Psi ; \Delta \To_i \Gamma$, then
    $t[\tau][\sigma] = t[\sigma][\tau]$.
  \item If $\ltyping i \delta {\Gamma'}$ and $\tau : \Psi ; \Delta \To_i \Gamma$, then
    $\delta[\tau][\sigma] = \delta[\sigma][\tau]$.
  \end{itemize}
\end{lemma}

Next, we move on to the global substitution lemma for terms and local substitutions. %
Prior to that, we must first show the local weakening lemma and local substitution lemma:
\begin{lemma}[Local Weakenings] $ $
  \begin{itemize}
  \item If $\ltyping i t T$ and $\tau : \Psi; \Delta \To_i \Gamma$, then
    $\ltyping[\Psi][\Delta] i{t[\tau]} T$. 
  \item If $\ltyping i \delta {\Delta'}$ and $\tau : \Psi; \Delta \To_i \Gamma$, then
    $\ltyping[\Psi][\Delta] i{\delta[\tau]}{\Delta'}$. 
  \end{itemize}
\end{lemma}

\begin{lemma}[Local Substitutions] $ $
  \begin{itemize}
  \item If $\ltyping i t T$ and $\ltyping[\Psi][\Delta] i \delta \Gamma$, then
    $\ltyping[\Psi][\Delta] i{t[\delta]} T$. 
  \item If $\ltyping i \delta \Delta$ and $\ltyping[\Psi][\Gamma'] i {\delta'} \Gamma$, then
    $\ltyping[\Psi][\Gamma'] i{\delta \circ \delta'} {\Delta}$. 
  \end{itemize}
\end{lemma}
\begin{proof}
  We do a mutual induction. %
  Notice that in this lemma, it is somewhat more cumbersome to establish the proof for
  local substitutions. %
  When $\delta = \cdot^m$, then we must reason about the properties of
  $\widecheck{\delta'}$. %
  If $\widecheck{\delta'} = g$ for some $g$, then we must show that both $\Gamma$ and
  $\Gamma'$ start with this $g$. %
  The details are given in the Agda mechanization. 
\end{proof}

Finally, we give the global substitution lemma:
\begin{lemma}[Global Substitutions] $ $
  \begin{itemize}
  \item If $\ltyping i t T$ and $\typing[\Psi']\sigma{\Psi}$, then
    $\ltyping[\Psi'][\Gamma[\sigma]] i{t[\sigma]}{T[\sigma]}$. 
  \item If $\ltyping i \delta \Delta$ and $\typing[\Psi']\sigma{\Psi}$, then
    $\ltyping[\Psi'][\Gamma[\sigma]] i{\delta[\sigma]}{\Delta[\sigma]}$. 
  \end{itemize}
\end{lemma}
\begin{proof}
  We do a mutual induction. %
  We consider a few interesting cases:
  \begin{itemize}[label=Case]
  \item
    \begin{mathpar}
      \inferrule
      {\ltyping 1 {s}{\cont[\Delta] T} \\
        \ljudge[\Psi] 1{T'} \\
        \ltyping[\Psi, u : (\judge[\Delta] T)][\Gamma[p(\id)]] 1 {t}{T'[p(\id)]}}
      {\ltyping 1 {\letbox u s t} T'}
    \end{mathpar}
    \begin{align*}
      & \ltyping[\Psi'][\Gamma[\sigma]] 1 {s[\sigma]}{\cont[\Delta] T [\sigma] = \cont[\Delta[\sigma]]{T[\sigma]} }
        \byIH \\
      & \typing[\Psi', u : (\judge[\Delta[\sigma]]{T[\sigma]})]{\sigma[p(\id)],
        u^{\id_{\Delta[\sigma[p(\id)]]}}/u}{\Psi, u : (\judge[\Delta] T)}
        \tag{by typing rule} \\
      & \ltyping[\Psi', u : (\judge[\Delta[\sigma]]{T[\sigma]})][\Gamma[p(\id)][\sigma[p(\id)],
        u^{\id_{\Delta[\sigma[p(\id)]]}}/u]] 1 {t[\sigma[p(\id)],
        u^{\id_{\Delta[\sigma[p(\id)]]}}/u]}{T'[p(\id)][\sigma[p(\id)],
        u^{\id_{\Delta[\sigma[p(\id)]]}}/u]}
        \byIH \\
      & \ltyping[\Psi', u : (\judge[\Delta[\sigma]]{T[\sigma]})][\Gamma[\sigma[p(\id)]]] 1 {t[\sigma[p(\id)],
        u^{\id_{\Delta[\sigma[p(\id)]]}}/u]}{T'[\sigma[p(\id)]]}
        \tag{by \Cref{lem:cvar:gsubst-skip}} \\
      & \ltyping[\Psi', u : (\judge[\Delta[\sigma]]{T[\sigma]})][\Gamma[\sigma][p(\id)]] 1 {t[\sigma[p(\id)],
        u^{\id_{\Delta[\sigma][p(\id)]}}/u]}{T'[\sigma][p(\id)]}
        \tag{by algebraic laws} \\
      & \ltyping[\Psi'][\Gamma[\sigma]] 1 {\letbox u s t}{T'[\sigma]}
        \tag{by typing rule}
    \end{align*}
    
  \item
    \begin{mathpar}
      \inferrule
      {\ljudge[\Psi] 1 \Gamma \\ \ltyping[\Psi, g: \Ctx][\Gamma[p(\id)]] 1 t T}
      {\ltyping{1}{\Lambda g. t}{(g : \Ctx) \STo T}}
    \end{mathpar}
    Basically we do the same as above, but instead of providing a global substitution
    of terms, we provide a global substitution of a local context.
    
  \item
    \begin{mathpar}
      \inferrule
      {\ltyping{1}{t}{(g : \Ctx) \STo T} \\ \ljudge[\Psi] 0 \Delta}
      {\ltyping{1}{t~\$~\Delta}{T[\id_{\Psi}, \Delta/g]}}
    \end{mathpar}
    \begin{align*}
      &\ltyping[\Psi'][\Gamma[\sigma]]{1}{t[\sigma]}{(g : \Ctx) \STo T[\sigma] = (g :
      \Ctx) \STo (T[\sigma[p(\id)], g/g])}
      \byIH \\
      & \ljudge[\Psi'] 0 {\Delta[\sigma]}
        \byIH \\
      & \ltyping[\Psi'][\Gamma[\sigma]]{1}{t~\$~\Delta[\sigma]}{T[\sigma[p(\id)],
        g/g][\id_{\Psi'}, \Delta[\sigma]/g]}
        \tag{by typing rule}
    \end{align*}
    Finally we must show the following equation:
    \begin{align*}
      T[\sigma[p(\id)],
      g/g][\id_{\Psi'}, \Delta[\sigma]/g] =
      T[\id_{\Psi}, \Delta/g][\sigma]
    \end{align*}
    We reason as follows:
    \begin{align*}
      T[\sigma[p(\id)], g/g][\id_{\Psi'}, \Delta[\sigma]/g]
      &= T[(\sigma[p(\id)], g/g) \circ (\id_{\Psi'}, \Delta[\sigma]/g)] \\
      &= T[(\sigma[p(\id)] \circ (\id_{\Psi'}, \Delta[\sigma]/g)), \Delta[\sigma] ] \\
      &= T[(\sigma \circ \id_{\Psi'}), \Delta[\sigma] ]
      \tag{by \Cref{lem:cvar:gsubst-skip}}\\
      &= T[\sigma, \Delta[\sigma] ]
        \tag{by right identity}
    \end{align*}
    On the right hand side,
    \begin{align*}
      T[\id_{\Psi}, \Delta/g][\sigma]
      &= T[(\id_{\Psi} \circ \sigma), \Delta[\sigma]/g] \\
      &= T[\sigma, \Delta[\sigma]/g]
        \tag{by left identity}
    \end{align*}
    Thus both sides agree. 
  \item
    \begin{mathpar}
      \inferrule
      {\ljudge i \Gamma \\ g : \Ctx \in \Psi \\ \text{$\Gamma$ ends with $g$} \\ |\Gamma| = m}
      {\ltyping i {\cdot_g^m}{\cdot}}
    \end{mathpar}
    In this case, we must look up $g$ in $\sigma$, and branch depending on
    $\sigma(g)$.
    \begin{itemize}[label=Subcase]
    \item
      If $\sigma(g) = g', \Delta$ meaning that $\sigma(g)$ ends with a context
      variable $g'$, then we must construct a typing judgment for $\cdot_{g'}^{m'}$
      for some $m'$. %
      In this case, say if $\Gamma = g, \Gamma'$, then $\Gamma[\sigma] = g', \Delta,
      (\Gamma'[\sigma])$. %
      Therefore, $m' = |\Delta| + |\Gamma'[\sigma]| = |\Delta| + m$. 
    \item If $\sigma(g) = \cdot, \Delta$, then we proceed similarly except that we
      must construct a typing judgment for $\cdot^{m'}$ instead. 
    \end{itemize}
    
  \end{itemize}
\end{proof}

\begin{lemma}[Distributivity of Global Substitutions] $ $
  \begin{itemize}
  \item If $\ltyping i t T$, $\ltyping[\Psi][\Delta] i \delta \Gamma$ and
    $\typing[\Phi] \sigma \Psi$, then
    $t[\delta][\sigma] = (t[\sigma][\delta[\sigma]])$.
  \item If $\ltyping i \delta \Delta$, $\ltyping[\Psi][\Gamma'] i {\delta'} \Gamma$
    and $\typing[\Phi] \sigma \Psi$, then
    $(\delta \circ \delta') [\sigma] = (\delta[\sigma]) \circ (\delta'[\sigma])$.
  \end{itemize}
\end{lemma}
\begin{proof}
  We do a mutual induction on $\ltyping i t T$ and $\ltyping i \delta \Delta$. %
  The difficulty is coming from the base cases of local substitutions. %
  We must incorporate the shape of $\delta'$ depending on the cases of $\delta$. %
  We refer the readers to the Agda mechanization for the detailed proof and how we do
  case analysis on the base cases. 
  
  % We consider these cases:
  % \begin{itemize}[label=Case]
  % \item
  %   \begin{mathpar}
  %     \inferrule
  %     {\ljudge i \Gamma \\ \text{$\Gamma$ ends with $\cdot$} \\ |\Gamma| = m}
  %     {\ltyping i {\cdot^m}{\cdot}}
  %   \end{mathpar}
  %   We have
  %   \begin{align*}
  %     (\delta \circ \delta') [\sigma]
  %     &= \cdot^{\widehat{\delta'}}_{\widecheck{\delta'}}[\sigma] \\
  %     (\delta[\sigma]) \circ (\delta'[\sigma])
  %     &= \cdot^m \circ (\delta'[\sigma]) \\
  %     &= \cdot^{\widehat{\delta'[\sigma]}}_{\widecheck{\delta'[\sigma]}}
  %   \end{align*}
  %   We do a further case analysis based on whether there is a $\widecheck{\delta'}$:
  %   \begin{itemize}[label=Subcase]
  %   \item $\widecheck{\delta'} = g$ for some $g$. %
  %     We need to further check whether $\sigma(g)$ ends with a $\cdot$ or another
  %     $g'$.
  %     \begin{itemize}[label=Subsubcase]
  %     \item $\sigma(g)$ ends with $\cdot$, then
  %       \begin{align*}
  %         \cdot^{\widehat{\delta'}}_{\widecheck{\delta'}}[\sigma]
  %         &= \cdot^{\widehat{\delta'} + |\sigma(g)|}
  %       \end{align*}
  %       Notice that $\widehat{\delta'} = \widehat{\delta'[\sigma]}$, we have the equation.
  %     \item $\sigma(g)$ ends with another $g'$, then
  %       \begin{align*}
  %         \cdot^{\widehat{\delta'}}_{\widecheck{\delta'}}[\sigma]
  %         &= \cdot^{\widehat{\delta'} + |\sigma(g)|}_{g'}
  %       \end{align*}
  %       Moreover, $\widecheck{\delta'[\sigma]} = g'$ as well by a simple induction
  %       argument, due to which we have the target goal.
  %     \end{itemize}
      
  %   \item $\widecheck{\delta'}$ does not exist. %
  %     Then
  %     \begin{align*}
  %       \cdot^{\widehat{\delta'}}_{\widecheck{\delta'}}[\sigma]
  %       &= \cdot^{\widehat{\delta'}} \\
  %       &= \cdot^{\widehat{\delta'[\sigma]}} \tag{by a simple induction} \\
  %       &= \cdot^{\widehat{\delta'[\sigma]}}_{\widecheck{\delta'[\sigma]}}
  %         \tag{$\widecheck{\delta'[\sigma]}$ also does not exist}
  %     \end{align*}
  %   \end{itemize}
    
  % \item
  %   \begin{mathpar}
  %     \inferrule
  %     {\ljudge i \Gamma \\ g : \Ctx \in \Psi \\ \text{$\Gamma$ ends with $g$} \\ |\Gamma| = m}
  %     {\ltyping i {\cdot_g^m}{\cdot}}
  %   \end{mathpar}
  %   \begin{align*}
  %     (\delta \circ \delta') [\sigma]
  %     &= \cdot^{\widehat{\delta'}}_{g}[\sigma] \\
  %     (\delta[\sigma]) \circ (\delta'[\sigma])
  %     &= (\cdot^m_g[\sigma]) \circ (\delta'[\sigma])
  %   \end{align*}
  %   We analyze what is the base case of $\sigma(g)$:
  %   \begin{itemize}[label=Subcase]
  %   \item $\sigma(g)$ ends with $\cdot$, then
  %     \begin{align*}
  %       \cdot^{\widehat{\delta'}}_{g}[\sigma]
  %       &= \cdot^{\widehat{\delta'}+|\sigma(g)|} \\
  %       (\cdot^m_g[\sigma]) \circ (\delta'[\sigma])
  %       &= \cdot^{m + |\sigma(g)|} \circ (\delta'[\sigma]) \\
  %       &= \cdot^{\widehat{\delta'[\sigma]}}_{\widecheck{\delta'[\sigma]}}
  %     \end{align*}
  %   \item
      
  %   \end{itemize}
  % \item
  %   \begin{mathpar}
  %     \inferrule
  %     {\ljudge i \Gamma \\ g : \Ctx \in \Psi \\ \text{$\Gamma$ ends with $g$} \\ |\Gamma| = m}
  %     {\ltyping i {\wk_g^m}{g}}
  %   \end{mathpar}
  % \end{itemize}
\end{proof}

\subsection{Equivalence Rules}

In this section, we describe the equivalence rules. %
They follow closely to the equivalence rules by \citet[Sec. 4]{hu2024layered}. %
We only show the rules for the newly added constructs.
\begin{mathpar}
  \inferrule
  {\ljudge[\Psi] 1 \Gamma \\ \ltyequiv[\Psi, g: \Ctx][\Gamma[p(\id)]] 1 t {t'} T}
  {\ltyequiv{1}{\Lambda g. t}{\Lambda g. t'}{(g : \Ctx) \STo T}}

  \inferrule
  {\ltyequiv{1}{t}{t'}{(g : \Ctx) \STo T} \\ \ljudge[\Psi] 0 \Delta}
  {\ltyequiv{1}{t~\$~\Delta}{t'~\$~\Delta}{T[\id_{\Psi}, \Delta/g]}}

  \inferrule
  {\ljudge[\Psi] 1 \Gamma \\ \ltyping[\Psi, g: \Ctx][\Gamma[p(\id)]] 1 t T \\ \ljudge[\Psi] 0 \Delta}
  {\ltyequiv{1}{(\Lambda g. t)~\$~\Delta}{t[\id_{\Psi}, \Delta/g]}{T[\id_{\Psi}, \Delta/g]}}

  \inferrule
  {\ltyping{1}{t}{(g : \Ctx) \STo T}}
  {\ltyequiv{1}{t}{\Lambda g. (t[p(\id)])~\$~g}{(g : \Ctx) \STo T}}
\end{mathpar}
In the rules above, we specify the congruence for meta-abstraction $\Lambda$ and the
meta-application. %
Moreover, they also have $\beta$ and $\eta$ rules in the expected ways.

We also have a equivalence judgment for local substitutions
\[
  \ltyequiv 1 \delta{\delta'} \Delta
\]
The only rules are the congruence rules determined by all possible constructors. %
Both equivalence judgments for terms and local substitutions must be mutually
defined. %

Note that there is no need to define the equivalence for global substitutions. %
Effectively, the equivalence for global substitutions is defined as the equality. %
This is because all terms stored in a global substitution are at layer $0$, and thus
they do not have meaningful dynamics. %
Therefore, equivalent global substitutions must also be equal. %

We first establish the presupposition lemma for equivalence:
\begin{lemma}[Presupposition] $ $
  \begin{itemize}
  \item If $\ltyequiv 1 t {t'} T$, then $\ltyping 1 t T$ and $\ltyping 1 {t'} T$.
  \item If $\ltyequiv 1 \delta {\delta'} \Delta$, then $\ltyping 1 \delta \Delta$ and $\ltyping 1 {\delta'} \Delta$.
  \end{itemize}
\end{lemma}
\begin{proof}
  We perform a mutual induction. %
  For the $\beta$ and $\eta$ rules, we apply substitution lemmas proved in the
  previous section. 
\end{proof}

\begin{lemma}[Local Weakenings] $ $
  \begin{itemize}
  \item If $\ltyequiv[\Psi][\Delta] 1 t{t'} T$ and $\tau : \Psi; \Gamma \To_1 \Delta$,
    then $\ltyequiv 1 {t[\tau]}{t'[\tau]} T$.
  \item If $\ltyequiv[\Psi][\Delta] 1 \delta{\delta'} \Delta'$ and
    $\tau : \Psi; \Gamma \To_1 \Delta$, then
    $\ltyequiv 1 {\delta[\tau]}{\delta'[\tau]} \Delta'$.
  \end{itemize}
\end{lemma}
\begin{proof}
  We follow the local weakening property for typing above. 
\end{proof}

A counterpart is w.r.t to global weakenings:
\begin{lemma}[Global Weakenings] $ $
  \begin{itemize}
  \item If $\ltyequiv[\Phi] 1 t{t'} T$ and $\gamma : \Psi \To_g \Phi$,
    then $\ltyequiv[\Psi][\Gamma[\gamma]] 1 {t[\gamma]}{t'[\gamma]}{T[\gamma]}$.
  \item If $\ltyequiv[\Phi] 1 \delta{\delta'} \Delta$ and $\gamma : \Psi \To_g \Phi$,
    then $\ltyequiv[\Psi][\Gamma[\gamma]] 1 {\delta[\gamma]}{\delta'[\gamma]}{\Delta[\gamma]}$.
  \end{itemize}
\end{lemma}
\begin{proof}
  Similarly, we follow the global weakening property for typing above. 
\end{proof}

\begin{lemma}[Congruence of Local Substitutions] $ $
  \begin{itemize}
  \item If $\ltyping 1 t T$ and $\ltyequiv[\Psi][\Delta] 1 \delta{\delta'} \Gamma$, then
    $\ltyequiv[\Psi][\Delta] 1{t[\delta]}{t[\delta']} T$. 
  \item If $\ltyping 1 {\delta} {\Delta'}$ and
    $\ltyequiv[\Psi][\Delta] 1 {\delta'}{\delta''} \Gamma$, then
    $\ltyequiv[\Psi][\Delta] 1{\delta \circ \delta'}{\delta \circ \delta''}
    {\Delta'}$.
  \end{itemize}
\end{lemma}
\begin{proof}
  We do a mutual induction on $\ltyping 1 t T$ and $\ltyping 1 {\delta} {\Delta'}$. %
  We only consider a few interesting cases.
  \begin{itemize}[label=Case]
  \item
    \begin{mathpar}
      \inferrule
      {\ltyping 1 {s}{\cont[\Delta'] T} \\
        \ljudge[\Psi] 1{T'} \\
        \ltyping[\Psi, u : (\judge[\Delta'] T)][\Gamma[p(\id)]] 1 {t}{T'[p(\id)]}}
      {\ltyping 1 {\letbox u s t} T'}
    \end{mathpar}
    In this case, we must use the global weakening lemma above to derive
    \[
      \ltyequiv[\Psi][\Delta[p(\id)]] 1 {\delta[p(\id)]}{\delta'[p(\id)]}{\Gamma[p(\id)]}
    \]
    and then we apply IH.
    
  \item
    For all base cases of local substitutions, we realize that given
    $\ltyequiv[\Psi][\Delta] 1 \delta{\delta'} \Gamma$,
    we have
    \begin{itemize}
    \item $\widecheck{\delta} = \widecheck{\delta'}$, and
    \item $\widehat\delta = \widehat{\delta'}$.
    \end{itemize}
    because they characterize $\Delta$, so the exact $\delta$ and $\delta'$ are
    irrelevant. 
    
    The target goal follows immediate. 
    
  \end{itemize}
\end{proof}

\begin{lemma}[Local Substitutions] $ $
  \begin{itemize}
  \item If $\ltyequiv 1 t {t'} T$ and $\ltyequiv[\Psi][\Delta] 1 \delta{\delta'} \Gamma$, then
    $\ltyequiv[\Psi][\Delta] 1{t[\delta]}{t'[\delta']} T$. 
  \item If $\ltyequiv 1 {\delta''} {\delta'''} {\Delta'}$ and
    $\ltyequiv[\Psi][\Delta] 1 \delta{\delta'} \Gamma$, then
    $\ltyequiv[\Psi][\Delta] 1{\delta'' \circ \delta}{\delta''' \circ \delta'}
    {\Delta'}$.
  \end{itemize}
\end{lemma}
\begin{proof}
  We proceed by a mutual induction on $\ltyequiv 1 t {t'} T$ and
  $\ltyequiv 1 {\delta''} {\delta'''} {\Delta'}$. %
  We only look into the $\beta$ and $\eta$ rule for meta-functions, because we cannot
  apply IH:
  \begin{itemize}[label=Case]
  \item
    \begin{mathpar}
      \inferrule
      {\ljudge[\Psi] 1 \Gamma \\ \ltyping[\Psi, g: \Ctx][\Gamma[p(\id)]] 1 t T \\ \ljudge[\Psi] 0 \Delta'}
      {\ltyequiv{1}{(\Lambda g. t)~\$~\Delta'}{t[\id_{\Psi}, \Delta'/g]}{T[\id_{\Psi}, \Delta'/g]}}
    \end{mathpar}
    \begin{align*}
      & \ltyequiv[\Psi, g: \Ctx][\Delta[p(\id)]] 1
        {t[\delta[p(\id)]]}{t[\delta'[p(\id)]]} T
        \tag{by local substitution lemma} \\
      & \ltyequiv[\Psi][\Delta]{1}{(\Lambda g. (t[\delta[p(\id)]]))~\$~\Delta'}{t[\delta'[p(\id)]][\id_{\Psi}, \Delta'/g]}{T[\id_{\Psi}, \Delta'/g]}
    \end{align*}
    Now we have to align up the right hand side. %
    The target right hand side is
    \[
      t[\id_{\Psi}, \Delta'/g][\delta']
    \]
    We reason as follows
    \begin{align*}
      t[\delta'[p(\id)]][\id_{\Psi}, \Delta'/g]
      &= t[\id_{\Psi}, \Delta'/g][\delta'[p(\id)][\id_{\Psi}, \Delta'/g]]
        \tag{by distributivity of global substitutions} \\
      &= t[\id_{\Psi}, \Delta'/g][\delta']
        \tag{by \Cref{lem:cvar:gsubst-skip}}
    \end{align*}
    and we have the target goal. 
  \item
    \begin{mathpar}
      \inferrule
      {\ltyping{1}{t}{(g : \Ctx) \STo T}}
      {\ltyequiv{1}{t}{\Lambda g. (t[p(\id)])~\$~g}{(g : \Ctx) \STo T}}
    \end{mathpar}
    We apply the local weakening lemma:
    \begin{align*}
      & \ltyequiv[\Psi][\Delta] 1 {t[\delta]}{t[\delta']} T 
    \end{align*}
    On the right hand side, we have
    \begin{align*}
      \Lambda g. (t[p(\id)])~\$~g[\delta']
      & = \Lambda g. (t[p(\id)][\delta'[p(\id)]])~\$~g \\
      & = \Lambda g. (t[\delta'][p(\id)])~\$~g
        \tag{by algebraic rule}
    \end{align*}
    
  \end{itemize}
\end{proof}

Next, we consider the global substitution lemma:
\begin{lemma}[Global Substitutions] $ $
  \begin{itemize}
  \item If $\ltyequiv 1 t {t'} T$ and $\typing[\Phi]\sigma\Psi$, then
    $\ltyequiv[\Phi][\Gamma[\sigma]] 1{t[\sigma]}{t'[\sigma]}{T[\sigma]}$. 
  \item If $\ltyequiv 1{\delta} {\delta'}{\Delta}$ and $\typing[\Phi]\sigma\Psi$, then
    $\ltyequiv[\Phi][\Gamma[\sigma]] 1{\delta[\sigma]}{\delta'[\sigma]}{\Delta[\sigma]}$. 
  \end{itemize}
\end{lemma}
\begin{proof}
  As we have noted above, there is no need for an equivalence judgment between two
  global substitutions. %
  Therefore we directly apply the same $\sigma$ on both sides. %
  Let us consider the $\beta$ and $\eta$ rules for meta-functions:
  \begin{itemize}[label=Case]
  \item
    \begin{mathpar}
      \inferrule
      {\ljudge[\Psi] 1 \Gamma \\ \ltyping[\Psi, g: \Ctx][\Gamma[p(\id)]] 1 t T \\ \ljudge[\Psi] 0 \Delta'}
      {\ltyequiv{1}{(\Lambda g. t)~\$~\Delta'}{t[\id_{\Psi}, \Delta'/g]}{T[\id_{\Psi}, \Delta'/g]}}
    \end{mathpar}
    \begin{align*}
      & \typing[\Phi, g : \Ctx]{\sigma[p(\id)], g/g}{\Psi, g : \Ctx} \\
      & \ltyping[\Phi, g: \Ctx][\Gamma[p(\id)][\sigma[p(\id)], g/g]] 1
        {t[\sigma[p(\id)], g/g]}{T[\sigma[p(\id)], g/g]}
        \tag{by global substitution lemma} \\
      & \ltyping[\Phi, g: \Ctx][\Gamma[\sigma][p(\id)]] 1
        {t[\sigma[p(\id)], g/g]}{T[\sigma[p(\id)], g/g]} \\
      &  \ljudge[\Phi] 0 {\Delta'[\sigma]}\\
      & \ltyequiv[\Phi][\Gamma[\sigma]]{1}{(\Lambda g. (t[\sigma[p(\id)], g/g]))~\$~(\Delta'[\sigma])}{t[\sigma[p(\id)], g/g][\id_{\Psi}, \Delta'[\sigma]/g]}{T[\sigma[p(\id)], g/g][\id_{\Psi}, \Delta'[\sigma]/g]}
    \end{align*}
    We now reason a few equations:
    \begin{align*}
      (\Lambda g. (t[\sigma[p(\id)], g/g]))~\$~(\Delta'[\sigma])
      = (\Lambda g. t)~\$~\Delta' [\sigma]
    \end{align*}
    and
    \begin{align*}
      t[\sigma[p(\id)], g/g][\id_{\Psi}, \Delta'[\sigma]/g]
      & = t[(\sigma[p(\id)], g/g) \circ (\id_{\Psi}, \Delta'[\sigma]/g)]  \\
      & = t[\sigma[p(\id)] \circ (\id_{\Psi}, \Delta'[\sigma]/g), \Delta'[\sigma]/g]  \\
      & = t[\sigma \circ \id_{\Psi}, \Delta'[\sigma]/g]  \\
      & = t[\id_{\Phi} \circ \sigma, \Delta'[\sigma]/g]  \\
      & = t[\id_{\Phi}, \Delta'/g][\sigma] 
    \end{align*}
    Similarly, we have
    \begin{align*}
      T[\sigma[p(\id)], g/g][\id_{\Psi}, \Delta'[\sigma]/g]
      = T[\id_{\Phi}, \Delta'/g][\sigma]
    \end{align*}
    This concludes the goal. 
  \item
    \begin{mathpar}
      \inferrule
      {\ltyping{1}{t}{(g : \Ctx) \STo T}}
      {\ltyequiv{1}{t}{\Lambda g. (t[p(\id)])~\$~g}{(g : \Ctx) \STo T}}
    \end{mathpar}
    \begin{align*}
      & \ltyping[\Phi][\Gamma[\sigma]]{1}{t[\sigma]}{(g : \Ctx) \STo T[\sigma]} \\
      & \ltyping[\Phi][\Gamma[\sigma]]{1}{t[\sigma]}{(g : \Ctx) \STo (T[\sigma[p(\id)],g/g])} \\
      & \ltyequiv[\Phi][\Gamma[\sigma]]{1}{t[\sigma]}{\Lambda g. t[\sigma][p(\id)]~\$~g}{(g : \Ctx) \STo (T[\sigma[p(\id)],g/g])} 
    \end{align*}
    We reason the right hand side
    \begin{align*}
      \Lambda g. (t[p(\id)])~\$~g[\sigma]
      &= \Lambda g. (t[p(\id)][\sigma[p(\id)], g/g])~\$~g \\
      &= \Lambda g. (t[\sigma[p(\id)]])~\$~g \\
      &= \Lambda g. (t[\sigma][p(\id)])~\$~g
    \end{align*}
    and hence establish the goal.
  \end{itemize}
\end{proof}

\subsection{Weak Head Reduction}

To compute the decide whether two terms are equivalence based on the typing rules
above, we take the approach of reducibility candidates. %
This approach requires a reduction strategy, and then we use a type directed
convertibility checking to determine whether two normal forms after reduction are
equivalent. %
In this approach, it is sufficient to have a term reduced to weak head normal forms,
specified as below:
\begin{alignat*}{2}
  w &:= &&\ v \sep \ze \sep \su t \sep \boxit t \sep \lambda x. t \sep \Lambda g. t
  \tag{Weak head normal form ($\Nf$)} \\
  v &:= &&\ x \sep u^\delta \sep v\ t \sep \letbox u v t \sep v~\$~\Gamma \tag{Neutral form
    ($\Ne$)}
\end{alignat*}
Basically the one-step weak head reduction simply takes the $\beta$ rules from the
equivalence relations and adds head reduction:

Head reductions:
\begin{mathpar}
  \inferrule*
  {\ltyred{t}{t'}{S \func T} \\ \ltyping 1{s}{S}}
  {\ltyred{t\ s}{t'\ s}{T}}

  \inferrule*
  {\ltyred {s}{s'}{\square T} \\ \ltyping[\Psi, u : T][\Gamma[p(\id)]] 1 {t}{T'[p(\id)]}}
  {\ltyred {\letbox u s t}{\letbox u{s'}{t}}{T'}}

  \inferrule*
  {\ltyred{t}{t'}{(g : \Ctx) \STo T} \\ \ljudge[\Psi] 0 \Delta}
  {\ltyred{t~\$~\Delta}{t'~\$~\Delta}{T[\id_{\Psi}, \Delta/g]}}
\end{mathpar}

$\beta$ reductions:
\begin{mathpar}
  \inferrule*
  {\ltyping[\Psi][\Gamma, x : S] 1 t T \\ \ltyping 1 s S}
  {\ltyred{(\lambda x. t)\ s}{t[\id_{\Gamma}, s/x]}{T}}

  \inferrule*
  {\ljudge[\Psi] 1 \Gamma \\ \ltyping[\Psi, g: \Ctx][\Gamma[p(\id)]] 1 t T \\ \ljudge[\Psi] 0 \Delta}
  {\ltyred{(\Lambda g. t)~\$~\Delta}{t[\id_{\Psi}, \Delta/g]}{T[\id_{\Psi}, \Delta/g]}}

  \inferrule*
  {\ltyping[\Psi][\cdot] 0 s T \\ \ljudge 1 \Gamma \\ \ljudge 1 T' \\ \ltyping[\Psi, u : T][\Gamma[p(\id)]] 1 {t}{T'[p(\id)]}}
  {\ltyred{\letbox u {\boxit s} t}{t[\id_{\Psi}, s/u]}{T'}}
\end{mathpar}

Notice that weak head reduction only occurs at layer $1$, so we do not need to rewrite
down the layer explicitly. %
This is because $\beta$ reduction only occurs at layer $1$ and all terms at layer $0$
are identified by their exact syntactic structure. %
We define the reflexive transitive closure $\ltyreds t{t'} T$ of weak head reduction
in the usual way. %
In an implementation, we would repeatedly do one-step weak head reduction, until it
reaches a normal form. %
To compare two terms of the same type, we would first compute the weak head normal
forms of both sides, and then based on their type, we perform a type-directed
convertibility check. %
This check also performs $\eta$ expansion when necessary, so the resulting algorithm
is complete w.r.t. the equivalence relation. %

Notice that the weak head reduction relation is a subrelation of the equivalence
above. %
Therefore, several theorems are simply carried over. %
We omit the proofs and only state the theorems below:
\begin{lemma}[Local Weakenings]
  If $\ltyred[\Psi][\Delta] t{t'} T$ and $\tau : \Psi; \Gamma \To_1 \Delta$,
    then $\ltyred {t[\tau]}{t'[\tau]} T$.
\end{lemma}

\begin{lemma}[Global Weakenings]
  If $\ltyred[\Phi] t{t'} T$ and $\gamma : \Psi \To_g \Phi$,
  then $\ltyred[\Psi][\Gamma[\gamma]] {t[\gamma]}{t'[\gamma]}{T[\gamma]}$.
\end{lemma}

\begin{theorem}[Preservation]
  If $\ltyred t {t'} T$, then $\ltyping 1 t T$ and $\ltyping 1 {t'} T$. 
\end{theorem}


\begin{lemma}[Local Substitutions]
  If $\ltyred t {t'} T$ and $\ltyping[\Psi][\Delta] 1 \delta \Gamma$, then
  $\ltyred[\Psi][\Delta] {t[\delta]}{t'[\delta]} T$. 
\end{lemma}
\begin{lemma}[Globale Substitutions]
  If $\ltyred t {t'} T$ and $\typing[\Phi]\sigma\Psi$, then
    $\ltyred[\Phi][\Gamma[\sigma]]{t[\sigma]}{t'[\sigma]}{T[\sigma]}$. 
\end{lemma}
\begin{lemma}[Uniqueness]
  If $\ltyred t {t'} T$ and $\ltyred t {t''} T$, then $t' = t''$. 
\end{lemma}
\begin{proof}
  Induction on $\ltyred t {t'} T$ and analyze $\ltyred t {t''} T$. 
\end{proof}


These conclusions can be generalized to its reflexive transitive closure. 

In the next section, we will first look into the logical relations which establish
normalization property, and then we show that this checking strategy indeed is
complete.
